\chapter{Roles 和 Tools}

\section{roles 的作用}

role 是可复用的 agent 角色包。它可以提供默认 agent 名、说明、provider 支持、工具需求和相关投影规则。用户可以安装 role，然后把 role 加到项目配置中。

\section{role 命令}

\begin{longtable}{p{0.36\linewidth}p{0.54\linewidth}}
\toprule
命令 & 作用 \\
\midrule
\cmd{ccb roles list} & 列出可用 roles。 \\
\cmd{ccb roles show <role_id>} & 显示 role 摘要。 \\
\cmd{ccb roles install [role_id]} & 安装 role。 \\
\cmd{ccb roles install --path PATH} & 从本地 path 安装。 \\
\cmd{ccb roles install --skip-tools} & 安装 role 但跳过 tools。 \\
\cmd{ccb roles update [role_id]} & 更新 role。 \\
\cmd{ccb roles sync [path]} & 从本地目录同步 roles。 \\
\cmd{ccb roles sync --with-tools} & 同步时包含 tools。 \\
\cmd{ccb roles doctor <role_id>} & 检查 role 和 tools 状态。 \\
\cmd{ccb roles add <role_spec>} & 把 role 写入当前项目配置。 \\
\cmd{ccb roles add <role_spec> --agent NAME --provider PROVIDER --window WINDOW} & 指定 agent 名、provider 或目标 window。 \\
\bottomrule
\end{longtable}

\section{Archi role 示例}

安装：

\begin{lstlisting}
ccb roles install agentroles.archi
\end{lstlisting}

加入项目：

\begin{lstlisting}
ccb roles add agentroles.archi:codex
\end{lstlisting}

这条命令会寻找最近的项目 \src{.ccb} anchor，并更新项目配置。也可以手动写 compact 配置：

\begin{lstlisting}
cmd | agentroles.archi:codex
\end{lstlisting}

加载时，CCB 会把 \cfg{agentroles.archi} 展开为 role 默认 agent name，并在 agent spec 中保留 role id。最终 daemon 按普通 agent 挂载它。

同一个 role 可以在一个项目里开多个本地 agent 实例。使用 \cmd{--agent NAME} 指定不同名字：

\begin{lstlisting}
ccb roles add agentroles.archi:codex --agent archi_review
ccb roles add agentroles.archi:claude --agent archi_qa
\end{lstlisting}

这种情况下，\cmd{ccb ask agentroles.archi ...} 会因为目标不唯一而失败；应直接 ask \src{archi_review} 或 \src{archi_qa}。

\section{role spec}

\cmd{roles add} 接受单个 role leaf。常见格式：

\begin{lstlisting}
agentroles.archi
agentroles.archi:codex
\end{lstlisting}

规则：

\begin{itemize}
  \item role id 应使用 publisher.role 形式。
  \item role spec 可以带 provider。
  \item role spec 不接受 workspace mode；workspace 需要在配置中单独设置。
  \item 可以用 \cmd{--agent NAME} 指定 agent 名；这是为同一个 role 创建第二个项目本地实例的正式方式。
  \item 可以用 \cmd{--provider PROVIDER} 覆盖 provider。
  \item 可以用 \cmd{--window WINDOW} 指定窗口。
\end{itemize}

\section{tools 命令}

当前 tools 命令支持 rich workbench：

\begin{longtable}{p{0.36\linewidth}p{0.54\linewidth}}
\toprule
命令 & 作用 \\
\midrule
\cmd{ccb update rich} & 安装、更新并启用 rich workbench。 \\
\cmd{ccb tools doctor workbench --profile rich} & 检查 rich workbench 状态。 \\
\cmd{ccb tools install workbench --profile rich} & 生成隔离的 WezTerm/Yazi/Markdown 配置。 \\
\cmd{ccb tools launch workbench --profile rich} & 启动 rich workbench。 \\
\bottomrule
\end{longtable}

常用环境变量：

\begin{itemize}
  \item \src{CCB_WORKBENCH_PROFILE}
  \item \src{CCB_WORKBENCH_FORCE_RICH}
  \item \src{CCB_WORKBENCH_THEME}
  \item \src{CCB_RICH_AUTO_START}
\end{itemize}

rich workbench 使用 XDG data/state/cache 目录下的 \src{ccb/tools/workbench} 维护状态。
